% sec-08: Document 08, the building code standards for the La Jolla facility.
\docpart{Building Code}{Building Code Standards for a PDAC LLM Oncology Clinical
Trial Facility}

\section{Scope, Occupancy, and Applicable Codes}

This document states what has to be built. It applies to a 12,400 gross square
foot leasehold with 7,850 square feet of clinical floor area, within an existing
permitted structure in the La Jolla Community Plan area.

\begin{mvtable}
\begin{tabularx}{\textwidth}{>{\raggedright\arraybackslash}p{4.2cm} >{\raggedright\arraybackslash}p{3.2cm} >{\raggedright\arraybackslash}X}
\headrow \thead{Area} & \thead{Occupancy} & \thead{Consequence of the classification}\\
\midrule
Robotic procedure suites, recovery, sterile core & Ambulatory care, Group B with ambulatory care provisions & Sprinklers throughout the smoke compartment; two means of egress from each suite; smoke compartment not exceeding 22,500 square feet \\
Exam rooms, infusion positions & Business, Group B & Standard egress; no additional smoke compartmentation \\
Investigational drug pharmacy & Business, Group B with hazardous materials control area & Containment and exhaust under USP 800; quantities below control area maximums \cite{usp797800} \\
Specimen laboratory & Business, Group B & Negative pressure; no flammable storage above control area limits \\
Machine room & Business, Group B & Dedicated cooling; restricted access under document 09 §2 \\
Administrative and staff & Business, Group B & Standard \\
\end{tabularx}
\end{mvtable}
\tabcapl{tab:occupancy}{Occupancy by area. The first row is the expensive one:
the ambulatory care classification, not the robots, drives the sprinkler,
egress, and smoke compartment cost in the tenant improvement budget at document
14 §6.}

The governing codes are the California Building Standards Code, Title 24
\cite{ccr24}; NFPA 99 for health care facilities \cite{nfpa99}; and NFPA 101 for
life safety \cite{nfpa101}. Where the California amendments and the model code
differ, the California amendment governs.

\section{Site and Structural Provisions}

\begin{mvtable}
\begin{tabularx}{\textwidth}{>{\raggedright\arraybackslash}p{4.4cm} >{\raggedright\arraybackslash}p{3.0cm} >{\raggedright\arraybackslash}X}
\headrow \thead{Element} & \thead{Requirement} & \thead{Why this trial needs it}\\
\midrule
Procedure suite live load & 100 pounds per square foot & Robot bases, boom-mounted equipment, and a full instrument cart on the same slab \\
Robot base anchorage & Post-installed anchors to a slab not less than 5 inches, designed for the manufacturer's overturning moment & Registration drift is a data quality failure before it is a safety failure \\
Vibration at the robot base & Not more than 3,200 micro-inches per second, one-third octave, during operation & Instrument tip deflection is the limiting factor for a needle-placement instance \\
Procedure suite clear height & 10 feet 6 inches to the lowest obstruction & Boom-mounted lights and the surgical platform's full articulation \\
Suite floor area & 720 square feet each, two suites & Envelope plus the 1.2 meter clearance in document 09 §3 plus a circulation path \\
Seismic restraint & All suspended equipment, medical gas, and the machine room racks & Title 24 seismic provisions; post-event inspection under document 11 §3 \\
\end{tabularx}
\end{mvtable}
\tabcapl{tab:structural}{Six structural provisions. The vibration limit is the
one a general contractor will question, and it is stated because the device
envelope work that this site inherits assumes it \cite{whipple60s,h2humanoid}.}

\section{Mechanical, Electrical, and Plumbing}

\begin{mvtable}
\begin{tabularx}{\textwidth}{>{\raggedright\arraybackslash}p{4.0cm} >{\raggedright\arraybackslash}p{1.9cm} >{\raggedright\arraybackslash}p{2.4cm} >{\raggedright\arraybackslash}X}
\headrow \thead{Space} & \thead{Air changes} & \thead{Pressure} & \thead{Additional requirement}\\
\midrule
Robotic procedure suite & 20 per hour & Positive, 0.01 inches water column & Minimum 4 outside air changes; temperature 68 to 75 degrees Fahrenheit \\
Sterile core & 20 per hour & Positive & Adjacent to both suites \\
Soiled utility & 10 per hour & Negative & Exhausted directly outdoors \\
Pharmacy buffer room & 30 per hour & Positive & ISO Class 7; USP 797 \cite{usp797800} \\
Pharmacy containment area & 12 per hour & Negative & External exhaust; USP 800 \cite{usp797800} \\
Specimen laboratory & 8 per hour & Negative & Fume hood exhausted directly outdoors \\
Machine room & 6 per hour & Neutral & 64 to 81 degrees Fahrenheit; dedicated cooling \\
\end{tabularx}
\end{mvtable}
\tabcapl{tab:mep}{Air changes and pressure by space. Document 09 §4 sets traffic
patterns against exactly these relationships, and document 13 §5 sets the
pharmacy and laboratory requirements against the same two rows.}

The essential electrical system carries three branches under NFPA 99: life
safety, critical, and equipment \cite{nfpa99}. Both procedure suites, the
machine room, the specimen freezers, and the pharmacy hoods are on the critical
or equipment branch. Isolated power is provided in each procedure suite. The
plumbing serving the pharmacy buffer and containment areas is separated from
general plumbing, with a dedicated hand-washing sink outside the buffer room
boundary.

\section{Medical Gas, Shielding, and Specialty Systems}

Each procedure suite is provided with oxygen, medical air, nitrogen for
instrument drive where the platform requires it, and two waste anesthetic gas
disposal connections. Vacuum is provided at four outlets per suite. Medical gas
systems are installed, tested, and verified under NFPA 99, and the verification
record is part of the commissioning package in §6 \cite{nfpa99}.

The protocol schedules cross-sectional imaging for response assessment under
RECIST 1.1 at intervals stated in document 13 §5 \cite{recist11}. That imaging is
performed off site under the imaging agreement, so this facility provides no
computed tomography or magnetic resonance suite and requires no structural
shielding for one. Intraoperative ultrasound is performed with a mobile imaging
assistant instance and requires no shielding. Stating this explicitly is
deliberate: a reader who assumes an imaging suite will over-build the mechanical
and electrical systems by a wide margin.

\section{Technology Infrastructure and the Machine Room}

The on-premises model node is why this document differs from an ordinary
ambulatory surgical center building code. The protocol requires the model to run
on premises rather than on a hosted service \cite{phase1protocol}.

\begin{mvtable}
\begin{tabularx}{\textwidth}{>{\raggedright\arraybackslash}p{4.0cm} >{\raggedright\arraybackslash}p{3.2cm} >{\raggedright\arraybackslash}X}
\headrow \thead{Provision} & \thead{Requirement} & \thead{Cross-reference}\\
\midrule
Room area & 180 square feet & Front matter site table \\
Power & Two 208 volt three-phase circuits on the equipment branch & NFPA 99 \cite{nfpa99} \\
Uninterruptible supply & 30 minutes at full load, then generator & Document 11 §4 \\
Cooling & Dedicated, redundant, 64 to 81 degrees Fahrenheit & §3 of this document \\
Physical access & Restricted zone; badge plus second factor & Document 09 §2 \\
Network segmentation & Three segments: clinical, model, robot control, with no route between model and robot control & Document 07 §5, and document 11 §5 \\
Egress of data & One monitored path to the audit trail store; no outbound internet from the model segment & Document 03 §5 \\
\end{tabularx}
\end{mvtable}
\tabcapl{tab:machineroom}{Seven machine room provisions. The sixth is the one
that makes the advisory-only boundary in document 07 §5 a matter of wiring
rather than of policy: there is no route from the model segment to the robot
control segment.}

\section{Commissioning and Certificate of Occupancy}

Commissioning proceeds as a gated sequence. The last commissioning gate is the
first activation gate in document 12 §1, so the two documents share a boundary
rather than overlapping.

\begin{mvtable}
\begin{tabularx}{\textwidth}{>{\raggedright\arraybackslash}p{0.9cm} >{\raggedright\arraybackslash}p{3.8cm} >{\raggedright\arraybackslash}X}
\headrow \thead{Gate} & \thead{Scope} & \thead{Evidence that closes it}\\
\midrule
C1 & Structural and anchorage & Special inspection reports; anchor pull test records \\
C2 & Mechanical air balance & Balance report showing every row of Table~\ref{tab:mep} met \\
C3 & Electrical and essential branch & Generator load bank test; transfer time record; isolated power line isolation monitor test \\
C4 & Medical gas & NFPA 99 verification certificate \cite{nfpa99} \\
C5 & Fire and life safety & Sprinkler acceptance; alarm test; egress inspection \cite{nfpa101} \\
C6 & Machine room and network & Segmentation test showing no route from model to robot control; cooling redundancy test \\
C7 & Robot installation and envelope survey & Envelope markings surveyed; interlock test record for all 14 instances \\
C8 & Certificate of occupancy & The certificate, issued after the county health plan check under document 05 §6 \\
\end{tabularx}
\end{mvtable}
\tabcapl{tab:commissioning}{Eight commissioning gates. Gate C6 is tested by
attempting a connection and failing to make one, which is the only form of
evidence that a segmentation claim can carry.}
